\section{A finite contraction with an explicit modulus}
\label{con:section}

This section uses the estimates proved above, with deliberately generous
constants.  Put
\[
 r(t)=e^{-t}\bigl(\psi(e^t)-e^t\bigr),\qquad S=2^{18}.
\]
The hypotheses needed here are
\begin{align}
 |r(t)|&\leq 4 &&(t\geq0),\label{con:global}\\
 \left|\int_u^v r(w)\,dw\right|&\leq20
       &&(0\leq u\leq v),\label{con:signed}\\
 |r(t)|&\leq \frac{2}{t^2}\int_0^t(t-u)|r(u)|\,du+\frac{S}{t}
       &&(t\geq2).\label{con:weighted}
\end{align}
There is also a useful one-sided regularity estimate.  Monotonicity of
$\psi$ gives, for $v\geq u$,
\[
 r(v)+1\geq e^{u-v}(r(u)+1).
\]
Since $0\leq\psi(e^u)e^{-u}\leq4$ and $1-e^{-h}\leq h$, it follows that
\begin{equation}
 r(v)\geq r(u)-5(v-u).\label{con:decline}
\end{equation}
The extra unit in the constant is harmless.  In particular, no two-sided
continuity assumption is being made: $r$ has upward jumps.
At rational spatial points $x\le y$, the same calculation is simply
$\psi(y)/y\ge\psi(x)/y$ together with $1-x/y\le\log(y/x)$;
no exponential evaluation is needed in its arithmetic implementation.

\subsection{A deficit detected by finitely many rational tests}

Fix an integer $q\geq1$, and define
\begin{equation}
 \varepsilon=\frac1q,\quad D=2^{20}q^3,\quad
 \delta=\frac1D,\quad \overline H=1024q,\quad
 Q=1+\frac1{320q},\quad m=327680q^2,\quad
 \sigma=\frac1{64q}.
 \label{con:parameters}
\end{equation}
Both $m$ and $m/2$ are integers.  Write $\lambda=\log Q$ and
$H=m\lambda$.  The elementary inequalities
\[
 \frac{z}{1+z}\leq\log(1+z)\leq z\qquad(z\geq0)
\]
imply
\begin{equation}
 \frac1{640q}\leq\lambda\leq\frac1{320q},
 \qquad 512q\leq H\leq1024q=\overline H.
 \label{con:block-length}
\end{equation}

\begin{lemma}[Finite block deficit]\label{con:deficit}
Suppose $\varepsilon\leq a\leq4$ and $|r(t)|\leq a$ throughout a
block $[s,s+H]$, where $e^s$ is rational.  Then
\begin{equation}
 \int_s^{s+H}\bigl(a-|r(t)|\bigr)\,dt
       \geq\frac1{64q^2}.
 \label{con:deficit-bound}
\end{equation}
The proof needs only finitely many comparisons of rational numbers.
\end{lemma}

\begin{proof}
The nodes
\[
 x_i=e^sQ^i,\qquad t_i=\log x_i=s+i\lambda,
       \qquad 0\leq i\leq m,
\]
have rational $x_i$.  Compute rational numbers $b_i$ satisfying
\begin{equation}
 |b_i-r(t_i)|\leq\sigma.\label{con:sample-error}
\end{equation}
Such numbers are obtained by a bounded prime-power sum and rational
logarithm enclosures.  Indeed, approximating each nonzero $\Lambda(n)$
with absolute error at most $\sigma$ produces an error at most
$\sigma\lfloor x_i\rfloor/x_i\leq\sigma$ after division by $x_i$.
All tests below concern the rational numbers $b_i$ and $a/4$.

Set $d(t)=a-|r(t)|$.  First suppose that
$b_i<a/4$ for all $0\leq i\leq m/2$.
For $t\in[t_i,t_{i+1}]$ in the first half of the block,
\eqref{con:decline}, applied from $t$ to $t_{i+1}$, gives
\[
 r(t)\leq r(t_{i+1})+5\lambda
       \leq a/4+\sigma+5\lambda
       \leq 9a/32<a/2.
\]
For $-a\leq z\leq a/2$ one has
$a-|z|\geq(a+z)/3$.  Therefore \eqref{con:signed} and
\eqref{con:block-length} imply
\[
 \int_s^{s+H/2}d(t)\,dt
 \geq\frac{aH/2-20}{3}
 \geq\frac{256-20}{3}
 \geq\frac1{64q^2}.
\]
This proves the claim in the first case.

Second, suppose that $b_i>-a/4$ for all $m/2\leq i\leq m$.
For each interval in the second half, apply \eqref{con:decline}
from its left endpoint to obtain
$r(t)\geq-a/4-\sigma-5\lambda\geq-9a/32>-a/2$.
The inequality $a-|z|\geq(a-z)/3$ for $-a/2\leq z\leq a$
gives the same lower bound for the integral of $d$ on this half.

If neither case applies, a bounded search finds an index $i_0$ in the
first half with $b_{i_0}\geq a/4$ and an index in the second half with
$b_i\leq-a/4$.  Let $v$ be the first index after $i_0$ with
$b_v\leq-a/4$, and let $u<v$ be the last index at or after $i_0$
with $b_u\geq a/4$.  These are minima and maxima of finite decidable
sets.  For $u<i<v$,
\[
 -a/4<b_i<a/4.
\]
For each $i$ from $u$ through $v-1$, the left sample is at least
$-a/4$ and the right sample is at most $a/4$.  Thus the two applications
of \eqref{con:decline} used in the preceding cases yield
\[
 |r(t)|\leq a/4+\sigma+5\lambda\leq9a/32<a/2
       \qquad(t_u\leq t\leq t_v).
\]
At the outer samples,
\[
 r(t_u)-r(t_v)\geq a/2-2\sigma\geq15a/32.
\]
Another application of \eqref{con:decline} gives
$t_v-t_u\geq15a/(32\cdot5)$.  Consequently
\[
 \int_{t_u}^{t_v}d(t)\,dt
 \geq\frac a2\frac{15a}{32\cdot5}
 =\frac{3a^2}{64}\geq\frac1{64q^2}.
\]
All three cases prove \eqref{con:deficit-bound}.  No sign decision for
an arbitrary real number and no intermediate-value principle occurs.
\end{proof}

\subsection{Using the deficit in the weighted inequality}

\begin{lemma}[One finite contraction]\label{con:one-step}
Let $s,t\geq0$ with $e^s$ and $e^t$ rational, and suppose
$\varepsilon\leq a\leq4$ and $|r(u)|\leq a$ for $s\leq u\leq t$.
If
\begin{equation}
 t\geq8(s+\overline H),\qquad
 t\geq D(8s+S),\qquad t\geq2,
 \label{con:step-threshold}
\end{equation}
then
\begin{equation}
 |r(t)|\leq a-3\delta.\label{con:step-margin}
\end{equation}
\end{lemma}

\begin{proof}
Place consecutive full blocks of length $H$ starting at $s$ and ending
no later than $t/2$.  Their number $b$ satisfies
\[
 b\geq \frac{t/2-s}{H}-1
   \geq\frac{3t}{8H}-1
   \geq\frac{t}{4H}
   \geq\frac{t}{4\overline H}.
\]
Here $s\leq t/8$ and $t\geq8\overline H\geq8H$ were used.
There is no need to decide a floor of a real number: when $x=e^t$ is
rational, the permitted blocks are found by the rational test
\[
 \bigl(e^s Q^{m i}\bigr)^2\leq x.
\]
The number of candidates can be bounded by $\lceil x\rceil+1$:
Bernoulli's finite inequality gives
$Q^{mi}\geq1+mi/(320q)=1+1024qi>i$.

On all these blocks the weight $t-u$ is at least $t/2$.
By Lemma~\ref{con:deficit},
\begin{align*}
 \int_0^t(t-u)|r(u)|\,du
 &\leq\frac{at^2}{2}+4st
      -b\frac{t}{2}\frac1{64q^2}\\
 &\leq\frac{at^2}{2}+4st
      -\frac{t^2}{8\overline H}\frac1{64q^2}.
\end{align*}
The term $4st$ bounds the possible excess above the baseline $a$
on $[0,s]$, using \eqref{con:global}.
Substitution in \eqref{con:weighted} gives
\[
 |r(t)|\leq a
    -\frac1{4\overline H\cdot64q^2}
    +\frac{8s+S}{t}
 =a-4\delta+\frac{8s+S}{t}
 \leq a-3\delta.
\]
This proves the stated, deliberately stronger, decrement.
\end{proof}

Every part of this argument refers only to arguments no larger than $x$.
This finite dependence is essential in the induction below.  The surplus
$2\delta$ in \eqref{con:step-margin}, beyond the decrement $\delta$
that will actually be used, also provides room for certified numerical
enclosures.  The rational-arithmetic implementation assigns an explicit
error budget to the compound expressions; the coarse sample precision
$\sigma$ in Lemma~\ref{con:deficit} is not the precision used for those
compound expressions.

\subsection{Integer thresholds and the modulus}

Set
\begin{equation}
 J=4D-D/q,\qquad a_j=4-j/D\quad(0\leq j\leq J),
 \label{con:levels}
\end{equation}
so that $J$ is an integer and $a_J=1/q$.  Define, by primitive recursion,
\begin{align}
 T_0&=2,\notag\\
 T_{j+1}&=16(T_j+\overline H)+D(16T_j+S)+1,\label{con:recurrence}\\
 P_j&=2^{2T_j}.\notag
\end{align}
These are integers.  Since $1/2\leq\log2\leq1$,
\begin{equation}
 T_j\leq\log P_j\leq2T_j.\label{con:dyadic-threshold}
\end{equation}
Starting at the rational point $P_j$, the hypotheses of
Lemma~\ref{con:one-step} are therefore met at every rational
$x\geq P_{j+1}$.  Indeed, putting $s=\log P_j$ and $t=\log x$,
\eqref{con:recurrence} bounds both $8(s+\overline H)$ and
$D(8s+S)$ from above, with a spare unit.
We retain only a decrement $\delta$ at each step.  After $J$ steps,
\begin{equation}
 |\psi(x)-x|\leq x/q\qquad(x\geq P_J).
 \label{con:psi-at-recursive-threshold}
\end{equation}
The arithmetic justification of this finite induction is given in the
next subsection.

The recurrence admits a simple closed upper bound.  Put $F=2^{20}D$.
From \eqref{con:recurrence} and $D=2^{20}q^3$ one checks
\[
 T_{j+1}+1\leq F(T_j+1),\qquad
 T_J+1\leq3F^J.
\]
For example, the coefficient of $T_j$ in the recurrence is
$16+16D<F$, and its constant term plus one is
$2^{14}q+2^{18}D+2<F$.
The elementary inequality $2^n\geq n^2$ for $n\geq4$ gives
$D\leq2^{D/16}$, because $D\geq2^{20}$.  Hence
\[
 F=2^{20}D\leq2^{D/8},\qquad
 3F^J\leq3\,2^{D^2/2}\leq2^{D^2},
\]
where $J\leq4D$ was used.  Thus $T_J\leq2^{D^2}$ and
\begin{equation}
 P_J\leq N_\psi(q):=2^{\,2^{D^2+1}},
       \qquad D=2^{20}q^3.
 \label{con:closed-modulus}
\end{equation}
In particular, the modulus is elementary, and therefore primitive
recursive.  No claim of numerical usefulness is intended.

\subsection{The induction predicate is a bounded finite check}
\label{con:bounded-induction}

Induction on the assertion ``for every $x\geq P_j$'' would not be an
instance of quantifier-free induction.  Instead, fix a final integer
ceiling $X$, as well as $q$, before starting the induction.  Form
rational lower and upper enclosures
$\underline\psi(n)\leq\psi(n)\leq\overline\psi(n)$
for $0\leq n\leq X$.  The stage predicate is the Boolean conjunction,
over the finite set of integers $P_j\leq n\leq X$, of
\begin{equation}
 \overline\psi(n)\leq(1+a_j)n,
 \qquad
 \underline\psi(n)\geq(1-a_j)(n+1).
 \label{con:endpoint-predicate}
\end{equation}
Denominators are cleared, so this is a primitive-recursive Boolean
function of its integer arguments.  Its equality to $1$ is a
quantifier-free formula.

The two checks enclose the value at the left of the cell $[n,n+1)$
and its limit at the right. Between consecutive integers, $\psi(x)$ is
constant and $\psi(x)/x-1$ decreases. When $a_j>1$ the lower bound
is also immediate from positivity. Therefore
\eqref{con:endpoint-predicate} controls every rational argument needed
in the next contraction whose value is less than $X+1$.
All mesh coordinates in Lemma~\ref{con:deficit} are rational, and their
indices and denominators have explicit primitive-recursive bounds.
An empty range in \eqref{con:endpoint-predicate} is checked as true.

The base envelope $a_0=4$ has slack: $0\leq\psi(x)\leq4x$ implies
$-1\leq r(t)\leq3$.  For the step, the analytic bound improves $a_j$
by $3\delta$, whereas the new check asks for improvement only by
$\delta$. In particular, each normalized logarithm enclosure may be
required to have error at most $\delta/8$. The lower end-of-cell
condition is tested at $z=n+1-\delta/64$, where $\psi(z)=\psi(n)$.
This point is at least $n$ and less than $X+1$, so both its threshold
and the finite ceiling are preserved. Section~\ref{sec:arithmetic}
gives the explicit algebra recovering the end-of-cell condition.
The remaining compound-expression errors are assigned a total budget
at most $\delta/4$ by the finite arithmetic compilation described
elsewhere in this report.  The surplus in
\eqref{con:step-margin} exceeds all these allowances.

Thus the stage implication is a finite arithmetic calculation with
parameters $q$ and $X$, and quantifier-free induction applies to its
Boolean check.  Only after the $J$ stages have been completed is the
arbitrary ceiling $X$ generalized.  This is the reason the proof does
not use an induction hypothesis with an unbounded universal quantifier.
